% ══════════════════════════════════════════════════════════════════════
% § 5  RG flow phases
% ══════════════════════════════════════════════════════════════════════

\subsection{Sign pattern of
  \texorpdfstring{$(\Bcond_{A},\Bcond_{B})$}{(B_A, B_B)}
  and the three topologies}
\label{sec:phases-topology}

At fixed UV matter content, an $\SU(N)_{0}\times\SU(N)_{1}$ theory that
is AF at both nodes ($b_{0}^{(0,1)}>0$) may realise one of three
distinct RG topologies in the coupling plane $(g_{0},g_{1})$, depending
on the signs of the boundary functions $\Bcond_{A}$ and $\Bcond_{B}$
defined in \eqref{eq:Bcond-def}:

\begin{enumerate}\setlength{\itemsep}{4pt}
  \item $(\Bcond_{A},\Bcond_{B})=(-,-)$: both axis single-node SCFTs
    are IR-unstable with respect to the orthogonal coupling.  The flow
    is released into the interior at each boundary and ends at a
    genuine two-coupling SCFT $C$.  This is \cite[Fig.~5]{IW2005},
    reproduced in Figure~\ref{fig:fig5}.
  \item Mixed ($\Bcond_{A}<0<\Bcond_{B}$ or the opposite):
    one boundary repels, the other traps.  The interior fixed point
    $C$ --- if the R-symmetry extremisation finds one --- is generically
    a saddle; a single axis SCFT is the global IR attractor.
  \item $(\Bcond_{A},\Bcond_{B})=(+,+)$: both axis SCFTs are IR-stable
    perturbations.  A separatrix through $C$ divides the basins of
    attraction for the two axis fixed points; $C$ itself is not an
    attractor, and generic flows in the interior end on one of the
    two axis SCFTs.  This is \cite[Fig.~6]{IW2005}, reproduced in
    Figure~\ref{fig:fig6}.
\end{enumerate}

The classification of the five interacting equal-rank non-Veneziano
classes (Table~\ref{tab:Bcond}) shows that class~$2$ realises
topology~(i) while classes $4,16,21,23$ all realise topology~(ii);
topology~(iii) is \emph{not realised} in the strict non-Veneziano
limit.

% Inline TikZ: Figure 5 topology (both axes IR-unstable, interior SCFT attracts)
\begin{figure}[h]
  \centering
  \begin{tikzpicture}[
    scale=1.4,
    flow/.style={thick, -{Stealth[length=3mm]},
      postaction={decorate},
      decoration={markings, mark=at position 0.55 with
        {\arrow{Stealth[length=2.5mm]}}}},
    fp/.style={circle, draw=black, fill=white, inner sep=1.6pt},
    fpIR/.style={circle, draw=red!70!black, fill=red!70!black, inner sep=2pt},
    fpUV/.style={circle, draw=blue!70!black, fill=blue!70!black, inner sep=2pt},
  ]
    % axes
    \draw[->, thick] (-0.3,0) -- (5.2,0) node[right] {$g_{1}$};
    \draw[->, thick] (0,-0.3) -- (0,4.5) node[above] {$g_{2}$};
    % fixed points
    \node[fpUV, label={below left:UV: $(0,0)$}] (UV) at (0,0) {};
    \node[fp,   label={left:$(0,\,g_{2}^{\star})$}]      (A1) at (0,3.2) {};
    \node[fp,   label={below:$(g_{1}^{\star},\,0)$}]     (A0) at (4.0,0) {};
    \node[fpIR, label={[align=left]above right:{\textbf{IR SCFT $C$}\\$(g_{1}^{\star\star},g_{2}^{\star\star})$}}]
          (IR) at (2.8,2.2) {};
    % flows from UV
    \draw[flow, blue!60!black] (UV) -- (A1);
    \draw[flow, blue!60!black] (UV) -- (A0);
    % off-axis flows away from boundary SCFTs
    \draw[flow, red!70!black] (A1) to[out=-20, in=165] (IR);
    \draw[flow, red!70!black] (A0) to[out=110, in=-75] (IR);
    % interior flow from UV
    \draw[flow, violet!70!black] (UV) to[out=50, in=225] (IR);
    % generic trajectories
    \draw[flow, gray!70, dashed] (0.6,0.2) to[out=60, in=210] (IR);
    \draw[flow, gray!70, dashed] (0.2,0.6) to[out=30, in=195] (IR);
    \draw[flow, gray!70, dashed] (1.2,0.1) to[out=80, in=240] (IR);
    \draw[flow, gray!70, dashed] (0.1,1.2) to[out=10, in=175] (IR);
    % labels
    \node[red!70!black]  at (0.95, 3.40) {\small $\alpha_{1}>0$};
    \node[red!70!black]  at (4.05, 0.48) {\small $\alpha_{2}>0$};
  \end{tikzpicture}
  \caption{Figure~5 topology: $(\alpha_{1},\alpha_{2})=(+,+)$, both
    axis SCFTs are IR-unstable and release the flow into the
    interior; the two-coupling SCFT $C$ at
    $(g_{1}^{\star\star},g_{2}^{\star\star})$ is the global IR
    attractor.  Realised by class~$2$ in the $\SU(N)\times\SU(N)$
    non-Veneziano landscape.}
  \label{fig:fig5}
\end{figure}


% Inline TikZ: Figure 6 topology (both axes IR-stable, separatrix through saddle C)
\begin{figure}[h]
  \centering
  \begin{tikzpicture}[
    scale=1.4,
    flow/.style={thick, -{Stealth[length=3mm]},
      postaction={decorate},
      decoration={markings, mark=at position 0.55 with
        {\arrow{Stealth[length=2.5mm]}}}},
    fp/.style={circle, draw=black, fill=white, inner sep=1.8pt},
    fpIR/.style={circle, draw=red!70!black, fill=red!70!black, inner sep=2pt},
    fpUV/.style={circle, draw=blue!70!black, fill=blue!70!black, inner sep=2pt},
  ]
    % axes
    \draw[->, thick] (-0.3,0) -- (5.4,0) node[right] {$g_{0}$};
    \draw[->, thick] (0,-0.3) -- (0,4.7) node[above] {$g_{1}$};
    % fixed points
    \node[fpUV, label={below left:UV: $(0,0)$}] (UV) at (0,0) {};
    \node[fpIR, label={[align=right]left:{\textbf{axis SCFT $A$}\\$(0,\,g_{1}^{\star})$}}]
          (A1) at (0,3.4) {};
    \node[fpIR, label={[align=center]below:{\textbf{axis SCFT $B$}\\$(g_{0}^{\star},\,0)$}}]
          (A0) at (4.2,0) {};
    \node[fp,   label={[align=left]above right:{\textbf{saddle $C$}\\(interior)}}]
          (C)  at (2.6,2.1) {};
    % flows from UV up the two axes
    \draw[flow, blue!60!black] (UV) -- (A1);
    \draw[flow, blue!60!black] (UV) -- (A0);
    % axis fixed points ATTRACT off-axis flows
    \draw[flow, red!70!black] (C) to[out=190, in=-10] (A1);
    \draw[flow, red!70!black] (C) to[out=-95, in=80] (A0);
    % incoming trajectories to C along stable manifold (separatrix)
    \draw[flow, violet!70!black, thick] (UV) to[out=42, in=225] (C);
    % separatrix (continues past C on the unstable direction)
    \draw[dashed, thick, violet!70!black]
      (0.3,0.8) to[out=-25, in=175] (C);
    \draw[dashed, thick, violet!70!black]
      (0.8,0.3) to[out=65, in=200] (C);
    % basin A trajectories (end on A1)
    \draw[flow, gray!70] (0.1,1.5) to[out=10, in=-15] (A1);
    \draw[flow, gray!70] (0.3,2.2) to[out=0, in=-40] (A1);
    \draw[flow, gray!70] (0.8,2.9) to[out=-10, in=-70] (A1);
    % basin B trajectories (end on A0)
    \draw[flow, gray!70] (1.5,0.1) to[out=80, in=105] (A0);
    \draw[flow, gray!70] (2.2,0.3) to[out=40, in=130] (A0);
    \draw[flow, gray!70] (2.9,0.8) to[out=20, in=160] (A0);
    % B label
    \node[red!70!black]  at (1.1, 3.60) {\small $\Bcond_{A}>0$};
    \node[red!70!black]  at (4.20, 0.50) {\small $\Bcond_{B}>0$};
  \end{tikzpicture}
  \caption{Topology (iii): $\Bcond_{A},\Bcond_{B}>0$.  Both boundary
    single-node SCFTs (filled red circles) are IR-stable; a separatrix
    (dashed violet) through the interior saddle $C$ divides the basins
    of attraction.  Realised by class~$16$ deformed by $\Nf$
    fundamentals on node~$1$ in the window
    $0.427\,N-1<\Nf<N-1$.}
  \label{fig:fig6}
\end{figure}


\subsection{Deforming class~$16$ by fundamentals on node~$1$}
\label{sec:phases-class16}

We now deform the representative theory of class~$16$ by adding
$\Nf$ pairs of (anti-)fundamentals $(\fund+\afund)^{\Nf}$ to node~$1$
only, keeping node~$0$'s matter unchanged.  The effect on the two
boundaries is asymmetric because $\Nf$ enters the decoupled-sub-theory
a-maximisation only when node~$1$ is the \emph{active} node:

\begin{itemize}\setlength{\itemsep}{4pt}
  \item \textbf{Boundary A ($g_{1}=0$, node~$0$ active).}  The
    boundary sub-theory consists of node~$0$'s own $A+2\bar A$ tensors
    and the three bifundamentals promoted to fundamentals of
    $\SU(N)_{0}$.  The N\_f fundamentals on node~$1$ are free spectators
    and do not participate in the boundary-A a-max.  The a-maximal
    bifundamental R-charge at boundary A is independent of $\Nf$:
    \begin{equation}
      R_{\text{bif}}^{(A)}\;=\;\tfrac{1+\sqrt{5}}{6}\;\simeq\;0.539\,.
    \end{equation}
    The corrected AF bound on the free node~$1$ is therefore
    \begin{equation}
      \Nf\;\le\;\alpha^{(A)}_{\Bcond}\,N
      +2\gamma^{(A)}_{\Bcond}\;=\;0.427\,N-1\,,
      \label{eq:class16-BcondA}
    \end{equation}
    where the coefficient $\alpha^{(A)}_{\Bcond}=3\cdot\tfrac{1}{2}
    \cdot 1\cdot(3R_{\text{bif}}^{(A)}-2)+\alpha^{(1)}_{0}$ with
    $\alpha^{(1)}_{0}=1$.  Thus
    \begin{equation}
      \Bcond_{A}<0\;\Longleftrightarrow\;\Nf<0.427\,N-1\,,
    \end{equation}
    i.e.\ exceeding the corrected AF bound flips boundary~A from
    IR-unstable to IR-stable.

  \item \textbf{Boundary B ($g_{0}=0$, node~$1$ active).}
    Node~$1$'s matter now includes its own $\bar A$ tensor, three
    bifundamentals (as flavours), \emph{and} the $\Nf$ fundamentals.
    The a-maximisation for this sub-theory changes as $\Nf$ is varied.
    At $\Nf=0$ the stored database value is
    $\alpha^{(B)}_{\Bcond}\,N+2\gamma^{(B)}_{\Bcond}=-0.573\,N+1$,
    corresponding to $\Bcond_{B}>0$ (axis~B IR-stable) at all $N\ge 2$.
    Adding $\Nf>0$ increases $R_{\text{bif}}^{(B)}$ towards the
    free-field value $2/3$, reducing $|3R_{\text{bif}}^{(B)}-2|$ and
    driving $\alpha^{(B)}_{\Bcond}$ towards $0$ from below;
    it does not change sign within the AF window $\Nf<N-1$.
\end{itemize}

The UV asymptotic-freedom of node~$1$ requires $\Nf<N-1$; node~$0$'s
AF is unchanged because none of its matter is altered.  Combining
these with the boundary analyses above yields the phase diagram
summarised in Table~\ref{tab:class16-phases}.

\begin{table}[h]
  \centering
  \begin{tabular}{lcccc}
    \toprule
    Window on $\Nf$ (node~$1$) & $b_{0}^{(0)}$ & $b_{0}^{(1)}$
                               & $\Bcond_{A}$  & $\Bcond_{B}$
                               \\
    \midrule
    $0\le \Nf<0.427\,N-1$ &$+$&$+$&$<0$&$>0$\\
    $0.427\,N-1<\Nf<N-1$ &$+$&$+$&$>0$&$>0$\\
    $\Nf\ge N-1$          &$+$&$\le 0$ &\multicolumn{2}{c}{--- not AF ---}\\
    \bottomrule
  \end{tabular}
  \caption{RG phase structure of class~$16$ as a function of the
    number of fundamental pairs $\Nf$ added to node~$1$.  The second
    row is the Figure~6 window.}
  \label{tab:class16-phases}
\end{table}

\subsection{Main result: explicit realisation of
  \texorpdfstring{\cite[Fig.~6]{IW2005}}{Intriligator--Wecht Figure 6}}
\label{sec:phases-main}

Combining the two boundary analyses with the UV-AF budget, we
conclude:

\begin{quote}
  \textbf{Theorem (main result).}  The $\SU(N)\times\SU(N)$
  quiver gauge theory with:
  \begin{itemize}\setlength{\itemsep}{1pt}
    \item node-$0$ matter $A+2\bar A$;
    \item node-$1$ matter $\bar A$ plus $\Nf$ vector-like pairs
      $(\fund+\afund)^{\Nf}$;
    \item three bifundamentals of structure
      $2\,(\fund_{0},\fund_{1})+(\afund_{0},\afund_{1})$;
  \end{itemize}
  is UV-asymptotically free and realises the two-boundary
  IR-stable RG topology of \cite[Fig.~6]{IW2005} in the window
  \begin{equation}
    0.427\,N-1\;<\;\Nf\;<\;N-1.
    \label{eq:main-theorem-window}
  \end{equation}
  For $N\ge 4$ this window is non-empty, and its width grows
  linearly with $N$ as $\approx 0.573\,N$.  In this window the
  two boundary fixed points attract generic flows; the interior
  R-symmetry extremum $C$ is not an attractor.
\end{quote}

This is the first explicit realisation, to our knowledge, of
\cite[Fig.~6]{IW2005}, answering the remark of those authors that
such examples had not been found.

\subsection{Remarks on the interior fixed point in the Figure~6 phase}
\label{sec:phases-interior}

The interior extremum $C$ exists algebraically throughout the window:
the symbolic $a$-maximisation of the deformed class~$16$ returns a
non-degenerate critical point in the R-charge polytope at every
$\Nf$ in \eqref{eq:main-theorem-window}.  However, the sign pattern
$(\Bcond_{A},\Bcond_{B})=(+,+)$ at both boundaries forces $C$ to be a
saddle of the two-coupling flow, \emph{not} a local attractor, and
hence $C$ does not describe the IR of generic trajectories in the
interior.  The two axis SCFTs --- the $\SU(N)_{0}$ single-node SCFT on
the $g_{0}$-axis and the $\SU(N)_{1}$ single-node SCFT on the
$g_{1}$-axis --- are both IR-stable attractors with finite basins
separated by the (co-dim $1$) stable manifold of $C$.

We emphasise that the Figure~6 phase is accessible within the
classified database by a \emph{single-parameter} deformation: only
$\Nf$ on node~$1$ is varied, and only one of the five interacting
classes (class~$16$) realises it.  The remaining four classes
($4$, $21$, $23$) in the mixed topology do not enter the Figure~6
window under analogous deformations, because either $\alpha^{(A)}_{\Bcond}$
or $\alpha^{(B)}_{\Bcond}$ has the wrong sign from the outset for the
corresponding deformation.

A fully quantitative characterisation of the separatrix basin
structure (position of $C$ in the $(g_{0},g_{1})$ plane, slope of the
separatrix near $C$, and basin area ratios) requires going beyond the
one-loop analysis employed here; this is left for future work.
